\SourceLeaf{345}
\par\medskip\Needspace{4\baselineskip}\pdfbookmark[2]{Art. 7. 8.}{nachlass-biquadratic-III-art-7-8}\centerline{\textbf{7. 8.}}\nobreak\smallskip
Diese Regel ist allgemein, welche Zahl $M$ auch bezeichnet. Für den Fall, der zunächst Gegenstand unserer Untersuchung sein soll, in dem $M$ als ungerade und von der Form $1+(2+2i)N$ vorausgesetzt wird, ist eine etwas abgeänderte Regel zweckmäßiger.

Man denke sich die Zahlen $f$ erneut in $4$ Klassen zerlegt. In die erste setzt man die $(h)$, deren Doppeltes ebenfalls in $f$ liegt; zur zweiten, $h'$, zählen wir diejenigen, deren Doppeltes $2h'$ zu $f'$ gehört. Ebenso bezeichnen $h''$ und $h'''$ diejenigen, deren Doppeltes zu $f''$ und $f'''$ gehört. Der Dezident $\varepsilon$ ist also
\[
\varepsilon=\Sigma\Theta\frac{2hM}{m}+\Sigma\Theta\frac{2h'M}{m}+\Sigma\Theta\frac{2h''M}{m}+\Sigma\Theta\frac{2h'''M}{m}
\]
Die Gesamtheit aller $2h$ und $-2h''+(1+i)m$ nennen wir $H$

die aller $-i(2h'-m)$ und $i(2h'''-im)$ nennen wir $H'$

$H$ und $H'$ umfassen somit alle $f$; jene sind die geraden, diese die ungeraden.

Ferner sind die folgenden Beziehungen anzuwenden
\[
\begin{aligned}
\Theta iN&=1+\Theta N\\
\Theta(-N)&=2+\Theta N\\
\Theta(-iN)&=3+\Theta N\\
\Theta(N+1)&=1-\Theta N\\
\Theta(N+1+i)&=2+\Theta N\\
\Theta(N+i)&=3-\Theta N
\end{aligned}
\]
folglich
\[
\begin{aligned}
\Theta\frac{(-2h'i+mi)M}{m}&=3-\Theta\frac{-2h'iM}{m}=-\Theta\frac{2h'M}{m}\\
\Theta\frac{(-2h''+m(1+i))M}{m}&=2+\Theta\frac{-2h''M}{m}=\Theta\frac{2h''M}{m}\\
\Theta\frac{(2h'''i+m)M}{m}&=1-\Theta\frac{2h'''iM}{m}=-\Theta\frac{2h'''M}{m}
\end{aligned}
\]
und
\[
\begin{aligned}
\varepsilon={}&\Sigma\Theta\frac{2hM}{m}-\Theta\frac{(-2h'i+mi)M}{m}\\
&+\Theta\frac{(-2h''+m(1+i))M}{m}-\Theta\frac{(2h'''i+m)M}{m}\\
={}&\Sigma\Theta\frac{HM}{m}-\Sigma\Theta\frac{H'M}{m}
\end{aligned}
\]

\SourceLeaf{346}
oder
\[
\varepsilon=\Sigma\pm\Theta\frac{fM}{m}
\]
wobei das obere Zeichen für gerade $f$, das untere für ungerade zu nehmen ist.

\Article{9}
Es sei nun allgemein $f=\xi+\eta i$. Die Zahlen $\xi,\eta$ sind durch die Bedingung, dass $f$ ein kleinster Rest von $m$ ist, oder dass nach dem Ansatz $\frac{f}{m}=x+yi$ die Größen $x$ und $y$ zwischen den Grenzen $0$ und $\frac12$ liegen müssen, auf bestimmte Bereiche beschränkt. Durch Unterscheidung der verschiedenen Fälle ließen sich dafür leicht bestimmte Regeln angeben. Geben wir $\eta$ einen bestimmten Wert, so hat wiederum $\xi$ seine bestimmten Grenzen. Nehmen wir zum Beispiel an, dass $a$ negativ und $b$ positiv ist, so müssen wegen
\[
\begin{aligned}
x&=\frac{a\xi+b\eta}{aa+bb}\\
y&=\frac{a\eta-b\xi}{aa+bb}
\end{aligned}
\]
die folgenden Bedingungen gelten:
\[
\begin{array}{rll}
\mathrm{I.}&\text{damit }x\text{ positiv werde}&\xi<-\frac{b}{a}\eta\\[3pt]
\mathrm{II.}&\text{damit }y\text{ positiv werde}&\xi<\frac{a}{b}\eta\\[3pt]
\mathrm{III.}&\text{damit }x<\frac12\text{ werde}&\xi>\frac{aa+bb-2b\eta}{2a}\\[3pt]
\mathrm{IV.}&\text{damit }y<\frac12\text{ werde}&\xi>\frac{2a\eta-aa-bb}{2b}
\end{array}
\]

Für positives $\eta$ schließt die zweite Bedingung bereits die erste ein, für negatives $\eta$ ist es umgekehrt. Ebenso ist die dritte Bedingung bereits in der vierten enthalten,
\[
\begin{array}{ll}
\text{wenn}&\eta<\tfrac12(a+b)\\
\text{und umgekehrt, wenn}&\eta>\tfrac12(a+b)
\end{array}
\]

Wir brauchen jedoch nicht alle acht Fälle, die hier auftreten können, einzeln zu betrachten. Für einen bestimmten Wert von $\eta$ bezeichnen wir lediglich die kleinere Grenze von $\xi$ mit $\xi^0$, die größere mit $\xi^{00}$, und bemerken nur, dass an diesen Grenzwerten stets entweder $x=0,y=0,x=\frac12,y=\frac12$ gilt, und zwar muss Folgendes gelten:
\SourceLeaf{347}
\[
\begin{array}{l|l|l}
\text{wenn}&\text{in der \emph{oberen} Grenze}&\text{in der \emph{unteren} Grenze}\\
a\text{ pos. }b\text{ positiv}&x=\frac12\text{ oder }y=0&x=0\text{ oder }y=\frac12\\
a\text{ neg. }b\text{ positiv}&x=0\text{ oder }y=0&x=\frac12\text{ oder }y=\frac12\\
a\text{ neg. }b\text{ negativ}&x=0\text{ oder }y=\frac12&x=\frac12\text{ oder }y=0\\
a\text{ pos. }b\text{ negativ}&x=\frac12\text{ oder }y=\frac12&x=0\text{ oder }y=0
\end{array}
\]
Der Kürze halber unterscheiden wir diese vier Fälle, indem wir sagen, dass im ersten $m$ zum ersten Quadranten gehört, im zweiten zum zweiten usw.

\Article{10}
Wir wollen nun die Summe aller $\pm\Theta\frac{fM}{m}$ näher betrachten, bei denen $\eta$ einen bestimmten Wert hat. Wenn $\xi$ nach und nach stetig vom kleinsten Wert $\xi^0$ zum größten $\xi^{00}$ wächst, ändert sich auch
\[
\frac{(\xi+\eta i)M}{m}=X+Yi
\]
nach dem Gesetz der Stetigkeit. Liegt $\frac{M}{m}$ im ersten Quadranten, so wachsen sowohl $X$ als auch $Y$ fortwährend; liegt $\frac{M}{m}$ im zweiten Quadranten, so nimmt $X$ fortwährend ab und $Y$ zu. Im dritten Quadranten tritt das Umgekehrte des ersten, im vierten das Umgekehrte des zweiten Falls ein. Die in $X+iY$ enthaltene ganze Zahl ändert sich jedoch sprunghaft, indem sich entweder $[X]$ oder $[Y]$ um eine Einheit ändert. Die Werte von $\xi$, bei denen ein solcher Übergang stattfindet, das heißt, bei denen entweder $X$ oder $Y$ eine ganze Zahl wird, seien der Reihe nach
\[
\xi',\ \xi'',\ \xi'''\dotsfour\xi^n
\]

Hier ist zu bemerken, dass weder diese Werte noch $\xi^0$ und $\xi^{00}$ ganze Zahlen sein können, außer für $\eta=0$, wo entweder $\xi^0$ oder $\xi^{00}=0$ wird. Es sei nun
\[
\begin{aligned}
\Theta\frac{(\xi'+\eta i)M}{m}-\Theta\frac{(\xi^0+\eta i)M}{m}&=\delta'\quad\text{(anders auszudrücken)}\\
\Theta\frac{(\xi''+\eta i)M}{m}-\Theta\frac{(\xi'+\eta i)M}{m}&=\delta''\\
&\text{usw.}\\
\Theta\frac{(\xi^{00}+\eta i)M}{m}-\Theta\frac{(\xi^{n-1}+\eta i)M}{m}&=\delta^n
\end{aligned}
\]
\SourceLeaf{348}
Dann sieht man leicht, da zwischen $\xi^0$ und $\xi'$ $[\frac12\xi']-[\frac12\xi^0]$ gerade und $[\frac12\xi'+\frac12]-[\frac12\xi^0+\frac12]$ ungerade ganze Zahlen liegen usw., dass bei Betrachtung allein des bestimmten Wertes von $\eta$ gilt:
\[
\begin{aligned}
(\pm1)\Sigma\Theta\frac{fM}{m}
={}&\{[\tfrac12\xi']-[\tfrac12\xi^0]-[\tfrac12\xi'+\tfrac12]+[\tfrac12\xi^0+\tfrac12]\}\\
&\qquad\pdot\Theta\frac{(\xi^0+\eta i)M}{m}\\
&+\{[\tfrac12\xi'']-[\tfrac12\xi']-[\tfrac12\xi''+\tfrac12]+[\tfrac12\xi'+\tfrac12]\}\\
&\qquad\pdot\{\Theta\frac{(\xi^0+\eta i)M}{m}+\delta'\}\\
&+\{[\tfrac12\xi''']-[\tfrac12\xi'']-[\tfrac12\xi'''+\tfrac12]+[\tfrac12\xi''+\tfrac12]\}\\
&\qquad\pdot\{\Theta\frac{(\xi^0+\eta i)M}{m}+\delta'+\delta''\}\\
&+\text{usw.}\\
&+\{[\tfrac12\xi^{00}]-[\tfrac12\xi^n]-[\tfrac12\xi^{00}+\tfrac12]+[\tfrac12\xi^n+\tfrac12]\}\\
&\qquad\pdot\{\Theta\frac{(\xi^0+\eta i)M}{m}+\delta'+\delta''+\dotstwo+\delta^n\}
\end{aligned}
\]
\[
\begin{aligned}
={}&-([\tfrac12\xi^0]-[\tfrac12\xi^0+\tfrac12])\pdot\Theta\frac{(\xi^0+\eta i)M}{m}\\
&-([\tfrac12\xi']-[\tfrac12\xi'+\tfrac12])\pdot\delta'\\
&-([\tfrac12\xi'']-[\tfrac12\xi''+\tfrac12])\pdot\delta''\\
&-\text{usw.}\\
&-([\tfrac12\xi^n]-[\tfrac12\xi^n+\tfrac12])\pdot\delta^n\\
&+([\tfrac12\xi^{00}]-[\tfrac12\xi^{00}+\tfrac12])\pdot\Theta\frac{(\xi^{00}+\eta i)M}{m}
\end{aligned}
\]
(Dabei gilt das obere Zeichen für gerades $\eta$, das untere für ungerades.)

\Needspace{17\baselineskip}Die Zahlen $\delta',\delta'',\delta'''$ usw. können keine anderen Werte als $+1$ und $-1$ haben. Den Wert $+1$ erhält $\delta'$ unter den folgenden Bedingungen, wobei die Werte von $X,Y$, die zu $\xi'$ gehören, mit $X',Y'$ bezeichnet sind:
\[
\begin{array}{l|l}
\frac{M}{m}\text{ im 1. Quadr.}&\frac{M}{m}\text{ im 2. Quadr.}\\
X'\text{ ganze gerade Zahl}&X'\text{ ganze gerade Zahl}\\
\text{und }[Y]\text{ ungerade}&[Y]\text{ gerade}\\\hline
Y'\text{ ganze gerade Zahl}&Y'\text{ ganze gerade Zahl}\\
\text{und }[X]\text{ gerade}&[X]\text{ gerade}
\end{array}
\]
\[
\begin{array}{l|l}
\frac{M}{m}\text{ im 3. Quadr.}&\frac{M}{m}\text{ im 4. Quadr.}\\
X'\text{ ganze gerade Zahl}&X'\text{ ganze gerade Zahl}\\
\text{und }[Y]\text{ gerade}&[Y]\text{ ungerade}\\\hline
Y'\text{ ganze gerade Zahl}&Y'\text{ ganze gerade Zahl}\\
\text{und }[X]\text{ ungerade}&[X]\text{ ungerade}
\end{array}
\]

\SourceLeaf{349}
Sobald sich eine dieser Bedingungen in die entgegengesetzte ändert, wird $\delta=-1$; ändern sich beide, so bleibt $\delta=+1$.

\Article{11}
Die folgenden Formeln dienen der bequemeren Übersicht dieser Rechnungen:
\[
\begin{aligned}
\text{es ist }m&=a+bi,\qquad aa+bb=d\\
M&=A+Bi,\qquad AA+BB=D\\
\frac{dM}{m}&=\alpha+\beta i,\qquad\alpha=aA+bB,\quad\beta=aB-bA\\
\frac{\xi+i\eta}{m}&=x+iy,\qquad M(x+iy)=X+iY
\end{aligned}
\]

Sind $\eta$ und $X$ gegeben, so gilt
\[
\begin{array}{rl}
1.&\xi=\frac{\beta\eta}{\alpha}+\frac{dX}{\alpha}\\[3pt]
2.&Y=\frac{D\eta}{\alpha}+\frac{\beta X}{\alpha}
\end{array}
\]
Sind $\eta$ und $Y$ gegeben, so gilt
\[
\begin{array}{rl}
3.&\xi=-\frac{\alpha\eta}{\beta}+\frac{dY}{\beta}\\[3pt]
4.&X=-\frac{D\eta}{\beta}+\frac{\alpha Y}{\beta}
\end{array}
\]
Sind $\eta$ und $x$ gegeben, so gilt
\[
\begin{array}{rl}
5.&\xi=-\frac{b\eta}{a}+\frac{dx}{a}\\[3pt]
6.&X=-\frac{B\eta}{a}+\frac{\alpha x}{a}\\[3pt]
7.&Y=\frac{A\eta}{a}+\frac{\beta x}{a}
\end{array}
\]
Sind $\eta$ und $y$ gegeben, so gilt
\[
\begin{array}{rl}
8.&\xi=\frac{a\eta}{b}-\frac{dy}{b}\\[3pt]
9.&X=\frac{A\eta}{b}-\frac{\alpha y}{b}\\[3pt]
10.&Y=\frac{B\eta}{b}-\frac{\beta y}{b}
\end{array}
\]

\Article{12}
Die Regel des Artikels $10$ lässt sich nun so ausdrücken. Erhält $\eta$ einen bestimmten Wert, so ist
\[
\Sigma\pm\Theta\frac{fM}{m}=k^0\Theta(X^0+Y^0i)-k^{00}\Theta(X^{00}+Y^{00}i)+\Sigma k
\]

\SourceLeaf{350}
Hier ist $k^0=0$, wenn $[\xi^0]$ gerade ist; $k^0=+1$, wenn $[\xi^0]$ ungerade und $\eta$ gerade ist; $k^0=-1$, wenn $[\xi^0]$ ungerade und $\eta$ ungerade ist. $k^{00}$ wird ebenso durch $[\xi^{00}]$ und $\eta$ bestimmt. Schließlich ist $\Sigma k$ eine Summe aus so vielen Zahlen, wie es zwischen $\xi=\xi^0$ und $\xi=\xi^{00}$ ganze Werte von $X$ oder $Y$ gibt. Jedes Mal ist $k=0$, wenn das entsprechende $[\xi]$ gerade ist, hingegen $=\pm1$, wenn $[\xi]$ ungerade ist. Das Zeichen wird wie folgt bestimmt. Ist $X$ eine ganze Zahl, so wird $k=1$, wenn zugleich gilt:
\[
\begin{array}{ll}
\eta&\text{gerade}\\
X&\text{gerade}\\
\relax[Y]&\text{gerade}\\
\frac{M}{m}&\text{im zweiten oder dritten Quadranten d.\,h. }\alpha\text{ negativ}
\end{array}
\]

Ist eine oder sind drei dieser Bedingungen nicht erfüllt, so wird $k=-1$; fehlen zwei oder alle vier, so bleibt $k=1$. Ist hingegen $Y$ eine ganze Zahl, so wird $k=1$, wenn von den $4$ Bedingungen
\[
\begin{array}{ll}
\eta&\text{gerade}\\
Y&\text{gerade}\\
\relax[X]&\text{gerade}\\
\frac{M}{m}&\text{im ersten oder zweiten Quadranten d.\,h. }\beta\text{ positiv}
\end{array}
\]
alle, zwei oder keine erfüllt sind.

\Article{13}
Jetzt müssen wir noch die Fälle einzeln betrachten, in denen $\xi^0$ oder $\xi^{00}$ (oder $X^0,Y^0,X^{00},Y^{00}$) eine ganze Zahl ist. Hier sind vier Fälle zu unterscheiden, wobei wir $a$ und $A$ als ungerade voraussetzen.

I. Liegt $m$ im ersten Quadranten, so ist für $x=\frac12,y=0$ die Größe $\eta=\frac12b$ eine ganze Zahl. Dann ist $Y^{00}=\frac12B$ eine ganze Zahl, und $\Theta(X^{00}+Y^{00}i)$ wird nur dann $=\Theta(\frac12A+\frac12Bi)$ sein, wenn $\beta$ negativ ist. Bei positivem $\beta$ wird hingegen $\Theta(\frac12A+(\frac12B-1)i)$ dafür genommen.

II. Liegt $m$ im zweiten Quadranten, so ergibt sich für $x=0,y=0$ der Wert $\eta=0$. Hier gilt für diesen Wert von $\eta$: $X^{00}=0,Y^{00}=0$. Man hat dann

\SourceLeaf{351}
\[
\begin{array}{rll}
\Theta(X^{00}+Y^{00}i)=&2,&\text{je nachdem }\dfrac{M}{m}\text{ im }1.\\
&3,&\hfill2.\\
&0&\hfill3.\\
&1&\hfill4.\text{ Quadr. liegt, und }k^{00}=1
\end{array}
\]

III. Liegt $m$ im dritten Quadranten, so ist für $x=\frac12,y=0$ die Größe $\eta=\frac12b$ eine ganze Zahl, wofür $X^0+Y^0i=\frac12A+\frac12Bi$ gilt. Man setzt dann
\[
\Theta(X^0+Y^0i)=\Theta(\tfrac12A+(\tfrac12B-1)i)
\]
sooft $\beta$ negativ ist.

IV. Liegt $m$ im vierten Quadranten, so ist für $\eta=0$
$\Theta(X^0+Y^0i=0,1,2,3$ zu setzen, je nachdem $\frac{M}{m}$ im $1.\ 2.\ 3.\ 4.$ Quadranten liegt;
$k^0=0$.

\Article{14}
Aus den vorhergehenden Untersuchungen folgt nun die folgende Bestimmung des Dezidenten.

Man sammle alle Werte von $x$ und $y$, die \emph{innerhalb} der Grenzen $0$ und $\frac12$ liegen und für die entweder $\eta$ und $X$ oder $\eta$ und $Y$ ganze Zahlen sind, und bestimme für jedes $x+iy$ nach den Regeln des Artikels $12$ den Wert von $k$.

Man sammle ferner alle Werte auf den Grenzen, das heißt, wo entweder $x=0$ oder $\frac12$, während $y$ \emph{zwischen} $0$ und $\frac12$ liegt, oder $y=0$ oder $=\frac12$, während $x$ \emph{zwischen} $0$ und $\frac12$ liegt, und die so beschaffen sind, dass $\eta$ eine ganze Zahl und $[\xi]$ ungerade ist. Das zugehörige $l$ bestimme man wie folgt. Es sei $\Theta M(x+yi)=\pm\theta$, das obere Zeichen für gerades, das untere für ungerades $\eta$
\[
\begin{array}{c|c|c|c|c}
&\multicolumn{4}{c}{\text{so ist für }m\text{ im}}\\
\text{für}&1.\text{ Quadr.}&2.\text{ Quadr.}&3.\text{ Quadr.}&4.\text{ Quadr.}\\
y=0&l=-\theta&l=-\theta&l=+\theta&l=+\theta\\
x=\frac12&l=-\theta&l=+\theta&l=+\theta&l=-\theta\\
y=\frac12&l=+\theta&l=+\theta&l=-\theta&l=-\theta\\
x=0&l=+\theta&l=-\theta&l=-\theta&l=+\theta
\end{array}
\]

\SourceLeaf{352}
Kürzer so:
\[
l=\pm\theta,
\]
\[
\begin{array}{ll}
\text{das Zeichen ist dasselbe wie das von }a&\text{wenn }x=0\\
\text{das entgegengesetzte}&\text{wenn }x=\frac12\\
\text{dasselbe wie das von }b&\text{wenn }y=\frac12\\
\text{das entgegengesetzte}&\text{wenn }y=0
\end{array}
\]

\Needspace{22\baselineskip}Zu $\Sigma k+\Sigma l$ kommt dann noch hinzu
\[
\left.\begin{array}{ll}
\text{wenn }m\text{ im zweiten Quadranten liegt:}&2,1,0,3\\
\text{wenn }m\text{ im vierten Quadranten liegt:}&0,
\end{array}\right\}
\begin{array}{l}\text{je nachdem }\frac{M}{m}\text{ im}\\1.\ 2.\ 3.\ 4.\text{ Quadr.}\end{array}
\]
wenn $m$ im ersten Quadranten liegt und $\frac12(a-1)$ ungerade ist,
\[
\begin{array}{ll}
\Theta(\frac12A+(\frac12B-1)i)&\text{wenn }\beta\text{ positiv}\\
\Theta(\frac12A+\frac12Bi)&\text{wenn }\beta\text{ negativ}
\end{array}
\]
wenn $m$ im dritten Quadranten liegt und $\frac12(a-1)$ ungerade ist,
\[
\begin{array}{ll}
-\Theta(\frac12A+\frac12Bi)&\text{wenn }\beta\text{ positiv}\\
-\Theta(\frac12A+(\frac12B-1)i)&\text{wenn }\beta\text{ negativ.}
\end{array}
\]
\BookRule{.3\textwidth}
